\newpage
\recipe{Neapolitan Pizza}{From \href{https://www.youtube.com/watch?v=xqVgJHThsco&t=195s}{Brian Lagerstrom's recipe}. Makes 4 pizzas. \underline{Dough needs to rise overnight!}}{}{}

\begin{ingredients}
	\item 450 g water
	\item 6 pinches active dry yeast (\sfrac{1}{8} tsp)
	\item 665 g `00' flour (or all purpose flour)
	\item 20 g salt
\end{ingredients}

\medskip

% --- Preparation ---
Add the \textbf{yeast} to the \textbf{water}\footnote{To properly activate the yeast, dissolve it in warm water with a teaspoon of sugar or flour and let it sit for 5--10 minutes, until foamy.}, then add the \textbf{flour} and \textbf{salt}.
Mix well with a spoon, then knead with a stand mixer\footnote{If you do not have a stand mixer, the dough can also be kneaded by hand. However, as with bread, the resulting crust will generally be less open, airy, and professional-looking.} until smooth and elastic.
Cover and let rest for 30 minutes.

Perform a strength-building fold to develop the gluten\footnote{\textit{Strength-building fold}: without removing the dough from its container, grab one side of the dough, stretch it upward, and fold it over to the opposite side. Repeat on all four sides, forming an ``X''. Turn the dough over so the seam is facing down before letting it rest again. Brian Lagerstrom has several videos, including the one linked under the title of this recipe, demonstrating this technique.}, then cover and let rest for another 30 minutes.

Perform another strength-building fold, then cover and let the dough proof at room temperature for 12 hours. 

Divide the dough into 4 balls and place them in lightly oiled containers in the refrigerator for at least 2 hours, or up to 5 days. 

Remove from the refrigerator 3 hours before making the pizza.
Stretch the pizzas on a well-floured surface, add the desired toppings\footnote{I recommend stretching the pizza (strictly by hand), then transferring it to a well-floured wooden or aluminum pizza peel before adding the toppings. Trying to lift a fully topped pizza directly from the work surface requires considerable skill and, in inexperienced hands, often ends in disaster. Also, try to leave the pizza on the peel for as little time as possible, especially after adding wet ingredients (such as tomato sauce), since the moisture will cause it to stick to the peel.}\footnote{Should the reader, or one of the prospective dinner guests, intend to top the pizza with \textit{pineapple}, I would like to point out that this is an utterly barbaric act, deserving of public humiliation.}, then bake in a preheated oven, ideally on a pizza stone, for about 10 minutes (the bottom should begin to develop nicely dark or slightly charred spots).

% --- Description ---
\begin{recipeblock}
\small	\noindent \textit{``In Naples they invented pizza. \\
		And everyone honks their horns, \\
		because they want to have pizza.'' \\ \normalfont \hspace{1.5cm} \textemdash the Rizzello Brothers.}
\end{recipeblock}